\chapter{Khảo sát và phân tích các giải pháp liên quan}
\label{ch:cong-trinh-lien-quan}

\section{Các hệ thống quản lý tài chính cá nhân hiện có và hạn chế}
Trong quá trình xây dựng giải pháp, nhóm tiến hành khảo sát các ứng dụng phổ biến trên thị trường, tiêu biểu như Money Lover, Spendee và Wallet by BudgetBakers:
\begin{itemize}
    \item \textbf{Money Lover (v4.0.x):} Ứng dụng nổi tiếng tại Việt Nam do ZooStudio phát triển. Hỗ trợ ghi chép chi tiêu, lập ngân sách định kỳ, quét hóa đơn qua ảnh và biểu đồ thống kê trực quan. Liên kết tài khoản ngân hàng chỉ có ở gói Premium.
    \item \textbf{Spendee (v5.1.x):} Giao diện hiện đại, tập trung vào ví dùng chung (Shared Wallets) cho gia đình hoặc nhóm bạn. Hầu hết tính năng nâng cao yêu cầu tài khoản trả phí.
    \item \textbf{Wallet by BudgetBakers (v8.4.x):} Cung cấp báo cáo tài chính chuyên sâu, hỗ trợ nhiều loại tiền tệ và đồng bộ hóa tự động với hơn 4.000 ngân hàng đối tác.
\end{itemize}

\begin{figure}[H]
    \centering
    \begin{minipage}{0.45\textwidth}
        \centering
        \includegraphics[height=8.5cm]{chapter2/money_lover_1.jpg}
        \centerline{(a) Giao diện ghi chép chi tiêu}
    \end{minipage}
    \hfill
    \begin{minipage}{0.45\textwidth}
        \centering
        \includegraphics[height=8.5cm]{chapter2/money_lover_2.jpg}
        \centerline{(b) Biểu đồ phân tích thu chi}
    \end{minipage}
    \caption{Giao diện khảo sát ứng dụng Money Lover}
    \label{fig:survey-money-lover}
\end{figure}

\begin{figure}[H]
    \centering
    \begin{minipage}{0.45\textwidth}
        \centering
        \includegraphics[height=8.5cm]{chapter2/spendee_1.jpg}
        \centerline{(a) Thiết lập hạn mức ngân sách}
    \end{minipage}
    \hfill
    \begin{minipage}{0.45\textwidth}
        \centering
        \includegraphics[height=8.5cm]{chapter2/spendee_2.jpg}
        \centerline{(b) Quản lý ví nhóm dùng chung}
    \end{minipage}
    \caption{Giao diện khảo sát ứng dụng Spendee}
    \label{fig:survey-spendee}
\end{figure}

\begin{figure}[H]
    \centering
    \begin{minipage}{0.31\textwidth}
        \centering
        \includegraphics[height=8.5cm]{chapter2/wallet_1.jpg}
        \centerline{(a) Tổng quan tài khoản}
    \end{minipage}
    \hfill
    \begin{minipage}{0.31\textwidth}
        \centering
        \includegraphics[height=8.5cm]{chapter2/wallet_2.jpg}
        \centerline{(b) Liên kết ngân hàng}
    \end{minipage}
    \hfill
    \begin{minipage}{0.31\textwidth}
        \centering
        \includegraphics[height=8.5cm]{chapter2/wallet_3.jpg}
        \centerline{(c) Biểu đồ báo cáo}
    \end{minipage}
    \caption{Giao diện khảo sát ứng dụng Wallet by BudgetBakers}
    \label{fig:survey-wallet}
\end{figure}
\subsection{Phương pháp và nghị trình khảo sát các giải pháp thị trường}
Nhóm thiết lập một quy trình khảo sát hệ thống (Survey Protocol) với các tiêu chí cụ thể:
\begin{itemize}
    \item \textbf{Phạm vi thời gian:} Khảo sát tiến hành từ tháng 09/2025 đến 02/2026. Ngày truy cập kiểm tra cuối cùng là 15/02/2026.
    \item \textbf{Từ khóa tìm kiếm:} Sử dụng từ khóa tiếng Anh và tiếng Việt trên Google Play, App Store và công cụ tìm kiếm: ``personal finance manager'', ``scholarship finder'', ``GPA calculator'', ``quản lý tài chính cá nhân'', ``học bổng sinh viên''.
    \item \textbf{Tiêu chí lựa chọn (Inclusion):}
    \begin{itemize}
        \item Phân hệ tài chính: Ứng dụng di động trên 100.000 lượt tải, hỗ trợ tiếng Anh hoặc tiếng Việt, tương thích thị trường Việt Nam.
        \item Phân hệ học bổng và học thuật: Cổng thông tin có lượng người dùng hoạt động thường xuyên hoặc là kênh chính thức của các cơ sở đại học lớn.
    \end{itemize}
    \item \textbf{Tiêu chí loại trừ (Exclusion):} Ứng dụng ngừng cập nhật trước năm 2024, giải pháp đóng phí bắt buộc 100\% ngay khi cài đặt, hoặc bảng tính Excel cá nhân không phân phối rộng rãi.
\end{itemize}

Để tổng hợp sự đối lập giữa các giải pháp hiện tại với Student360, nhóm xây dựng bảng ma trận so sánh tính năng cốt lõi trong Bảng~\ref{tab:so-sanh-giai-phap}.

\begin{table}[H]
\centering
\begin{tabular}{|L{3.2cm}|C{2.0cm}|C{2.0cm}|C{2.0cm}|L{2.5cm}|L{2.6cm}|}
\hline
\textbf{Tiêu chí đối sánh} & \textbf{Money Lover} & \textbf{Spendee} & \textbf{Wallet} & \textbf{Các giải pháp truyền thống} & \textbf{Đề tài Student360} \\ \hline
Ghi chép giao dịch & Có (Miễn phí / Trả phí) & Có (Miễn phí / Trả phí) & Có (Miễn phí / Trả phí) & Không hỗ trợ & Có (Miễn phí hoàn toàn) \\ \hline
Tự động hóa chia 6 lọ tài chính & Một phần (phải cấu hình ví thủ công) & Không hỗ trợ & Không hỗ trợ & Không hỗ trợ & Có (Hỗ trợ cấu hình và chia tự động) \\ \hline
Phân loại bằng AI & Thô sơ (Regex / Khớp chuỗi) & Không hỗ trợ & Không hỗ trợ & Không hỗ trợ & Có (LLM Few-shot + Rules) \\ \hline
Cảnh báo ngân sách chủ động & Có (Hạn mức tĩnh theo tháng) & Có (Hạn mức theo ví) & Có (Thống kê tĩnh) & Không hỗ trợ & Có (Đa lớp, thời gian thực qua FCM) \\ \hline
Liên kết học bổng & Không hỗ trợ & Không hỗ trợ & Không hỗ trợ & Lọc cứng thủ công qua tin bài & Có (So khớp hồ sơ tự động) \\ \hline
Quản lý học thuật & Không hỗ trợ & Không hỗ trợ & Không hỗ trợ & Tách biệt trên cổng Student Portal & Có (Tính toán GPA và tiến độ tích lũy) \\ \hline
Tương tác AI Agent & Không hỗ trợ & Không hỗ trợ & Không hỗ trợ & Không hỗ trợ & Có (Trò chuyện thời gian thực qua SSE) \\ \hline
\end{tabular}

\caption{Bảng ma trận đối sánh tính năng giữa các giải pháp thị trường và đề tài Student360}
\label{tab:so-sanh-giai-phap}
\end{table}

\subsection{Hạn chế đối với đối tượng sinh viên}
Các ứng dụng trên đã giải quyết được nhiều nhu cầu cơ bản, nhưng vẫn tồn tại các hạn chế cốt lõi đối với sinh viên:
\begin{enumerate}
    \item \textbf{Thiếu tích hợp sâu mô hình 6 lọ:} Người dùng phải tự tính toán tỷ lệ phân bổ và chia nhỏ dòng tiền thủ công mỗi khi có thu nhập mới. Chưa ứng dụng nào số hóa quy chuẩn mô hình ``6 cái lọ'' của T. Harv Eker \cite{eker2005-6jars}.
    \item \textbf{Mức độ tự động hóa bằng AI còn thấp:} Đa số chỉ dùng thuật toán so khớp chuỗi (string matching) hoặc biểu thức chính quy để phân loại giao dịch. Khi người dùng nhập mô tả viết tắt hoặc sai chính tả, hệ thống không nhận diện được và buộc chọn danh mục thủ công.
    \item \textbf{Thiếu tính chủ động và cá nhân hóa (Proactivity):} Các ứng dụng hoạt động theo mô hình thụ động -- chỉ ghi nhận và thống kê sau khi giao dịch xảy ra. Hệ thống thiếu cơ chế phân tích hành vi ngầm để cảnh báo chi tiêu vượt hạn mức hoặc gợi ý tái cấu trúc dòng tiền trước khi sự cố xảy ra.
\end{enumerate}

\section{Các giải pháp kết nối học bổng truyền thống}
Đối với sinh viên Việt Nam, thông tin học bổng thường được công bố qua các kênh truyền thống:
\begin{itemize}
    \item Trang web chính thức của Phòng Công tác Sinh viên (CTSV) hoặc Giáo vụ khoa của các trường Đại học.
    \item Các diễn đàn trực tuyến, mạng xã hội (Facebook, LinkedIn).
    \item Các trang tin học bổng chuyên biệt (ví dụ: Scholarship Planet, YBOX).
\end{itemize}

Các phương thức này bộc lộ nhiều hạn chế:
\begin{enumerate}
    \item \textbf{Thông tin bị phân mảnh và chậm trễ:} Sinh viên phải truy cập thủ công nhiều trang web. Nhiều học bổng có thời gian ứng tuyển ngắn, sinh viên thường chỉ phát hiện khi sát hạn nộp.
    \item \textbf{Cơ chế lọc thô sơ:} Các trang tin chỉ hỗ trợ lọc theo tiêu chí cứng cơ bản (Quốc gia, Ngành học chung). Sinh viên không biết hồ sơ của mình đáp ứng bao nhiêu phần trăm điều kiện của học bổng.
    \item \textbf{Gánh nặng cho nhà tài trợ và ban quản trị:} Quy trình xét duyệt hồ sơ đa phần vẫn thực hiện thủ công qua email hoặc Google Forms, gây khó khăn quản lý tập trung và dễ xảy ra sai sót.
\end{enumerate}

\section{Các giải pháp quản lý học thuật và bảng điểm hiện có}
Các giải pháp hiện có cho nhu cầu theo dõi kết quả học tập thuộc một trong hai nhóm:
\begin{itemize}
    \item \textbf{Cổng thông tin sinh viên (Student Portal):} Kênh chính thức để tra cứu điểm số sau khi nhà trường công bố. Đảm bảo tính chính xác nhưng không hỗ trợ sinh viên tự nhập điểm dự kiến để mô phỏng kịch bản GPA, và hoàn toàn tách biệt với nhu cầu tài chính hay học bổng.
    \item \textbf{Ứng dụng tính GPA độc lập:} Cho phép nhập điểm và số tín chỉ để tính GPA tích lũy. Tuy nhiên không đồng bộ với danh mục môn học chuẩn hóa của từng trường, dễ dẫn đến sai lệch khi trường có môn trùng tên nhưng khác số tín chỉ.
\end{itemize}

Cả hai nhóm đều tồn tại hạn chế chung:
\begin{enumerate}
    \item \textbf{Thiếu liên kết giữa dữ liệu học thuật và các nhu cầu khác:} GPA là tiêu chí xét duyệt quan trọng của nhiều học bổng, nhưng sinh viên phải tự tổng hợp thủ công từ cổng trường để điền vào hồ sơ ứng tuyển ở hệ thống khác -- không có cơ chế đối chiếu tự động.
    \item \textbf{Không chuẩn hóa dữ liệu môn học giữa các trường:} Do mỗi trường có mã môn, số tín chỉ và thang điểm riêng, các ứng dụng tính GPA độc lập không thể tự động xác thực dữ liệu, dẫn đến sai số tích lũy.
\end{enumerate}

\section{Nghiên cứu liên quan và ứng dụng trí tuệ nhân tạo}
\label{sec:nghien-cuu-lien-quan}
Sự bùng nổ của Mô hình ngôn ngữ lớn (LLM) và hệ thống đại lý thông minh (AI Agents) đã mở ra hướng giải quyết mới cho bài toán quản lý và gợi ý thông tin trong giáo dục. Phần này tổng hợp các nghiên cứu liên quan gần đây theo nhóm phương pháp và kết quả.

\subsection{Lĩnh vực quản lý tài chính cá nhân (Personal Finance Management)}
\begin{itemize}
    \item \textbf{Phương pháp:} Nghiên cứu truyền thống áp dụng SVM, Naive Bayes hoặc cây quyết định để phân loại văn bản giao dịch. Xu hướng gần đây chuyển sang các mô hình Transformer \cite{vaswani2017attention} như BERT \cite{devlin2018bert} hoặc LLM thông qua kỹ thuật Few-shot Prompting \cite{brown2020-gpt3}.
    \item \textbf{Dữ liệu:} Chuỗi văn bản mô tả giao dịch ngắn (dưới 15 từ) từ người dùng di động, chứa nhiều từ viết tắt hoặc lỗi chính tả (ví dụ: ``mua banh mi 15k'', ``tna 50k'').
    \item \textbf{Kết quả:} LLM kết hợp Few-shot Prompting và Chain-of-Thought \cite{wei2022-cot} cải thiện F1-Score từ 15\% đến 25\% so với phương pháp cũ mà không cần fine-tuning lại mô hình.
\end{itemize}

\subsection{Lĩnh vực gợi ý học bổng (Scholarship Recommender Systems)}
\begin{itemize}
    \item \textbf{Phương pháp:} Hai hướng chính: Lọc dựa trên nội dung (Content-based filtering) dùng TF-IDF, Cosine hoặc Jaccard \cite{jaccard1901}, và Lọc cộng tác (Collaborative filtering). Do dữ liệu tương tác thưa thớt, các nghiên cứu hiện đại ưu tiên hệ gợi ý kết hợp (Hybrid Recommender Systems) \cite{burke2002hybrid, adomavicius2005toward} để khắc phục Cold Start.
    \item \textbf{Kết quả:} Kết hợp Fuzzy Matching cấp độ token (Jaccard) và ký tự (N-gram \cite{cavnar1994-ngram}, Levenshtein \cite{levenstein1966}) giúp hệ thống chịu lỗi viết tắt tốt hơn. Các hệ thống hybrid đạt Precision@5 trên 80\% trong môi trường thử nghiệm học thuật.
\end{itemize}

\subsection{Lĩnh vực Tác nhân AI hỗ trợ (LLM-based Agents)}
\begin{itemize}
    \item \textbf{Phương pháp:} Kiến trúc Tác nhân AI tự trị (Autonomous Agents) \cite{wang2023survey-agent} theo mô hình ReAct (Reasoning and Acting) \cite{yao2022-react}: LLM thực hiện vòng lặp suy luận ý định $\rightarrow$ gọi công cụ bên ngoài (truy vấn CSDL, gọi API) $\rightarrow$ nhận kết quả $\rightarrow$ tiếp tục suy luận hoặc đưa ra câu trả lời.
    \item \textbf{Kết quả:} ReAct kết hợp Tool-calling giúp AI Agent truy vấn số dư thực tế, đưa ra nhận định tài chính cá nhân hóa và định tuyến chính xác ý định tìm kiếm học bổng.
\end{itemize}

\section{Các giải pháp thông báo đẩy thời gian thực trên nền tảng di động}
Hệ thống thông báo đẩy (Push Notification) đóng vai trò quan trọng trong việc tăng tỷ lệ giữ chân người dùng. Các công nghệ được Student360 áp dụng:
\begin{itemize}
    \item \textbf{Firebase Cloud Messaging (FCM) \cite{fcm-docs}:} Dịch vụ đám mây của Google, gửi thông báo đến thiết bị Android và iOS kể cả khi ứng dụng đã bị đóng (Background/Killed state). FCM tối ưu hóa thời lượng pin bằng cách gom nhóm kết nối mạng ngầm. Đây là lựa chọn chính để gửi thông báo học bổng và cảnh báo tài chính.
    \item \textbf{BullMQ và Redis Queue \cite{bullmq-docs}:} Để xử lý gửi thông báo đến hàng nghìn thiết bị mà không làm nghẽn tiến trình chính, hệ thống áp dụng kiến trúc hàng đợi thông điệp. BullMQ kết hợp Redis \cite{redis-docs} quản lý thứ tự gửi, xử lý lỗi tự động và phân bổ tài nguyên tối ưu. Điểm khác biệt cốt lõi là kiến trúc thông báo \textbf{chủ động đa lớp} với nhiều loại thông báo Cron tự động:
    \begin{enumerate}
        \item \textbf{Cảnh báo ngưỡng số dư hũ} (theo số tiền tuyệt đối hoặc phần trăm ngân sách): kích hoạt ngay sau mỗi giao dịch chi tiêu theo kiểu fire-and-forget.
        \item \textbf{Cảnh báo ngân sách hũ sắp hết hoặc đã cạn}: gửi khi số dư còn lại $\leq 10\%$ hoặc đã âm.
        \item \textbf{Nhắc lịch chuyển tiền tự động sắp đến}: gửi 24 giờ trước ngày thực hiện, kèm xác nhận sau khi chuyển thành công hoặc thất bại.
        \item \textbf{Tóm tắt chi tiêu tuần} (Cron thứ Hai 8:00): tổng hợp tổng thu, tổng chi và hũ chi nhiều nhất trong tuần.
        \item \textbf{Báo cáo tài chính cuối tháng} (Cron ngày 1 hàng tháng 8:00): nhắc xem phân tích AI 6 lọ.
        \item \textbf{Nhắc nhở cập nhật tài chính} (Cron thứ Hai 9:00): gửi cho người dùng không nhập giao dịch nào trong 7 ngày gần nhất.
        \item \textbf{Nhắc lịch chuyển tiền sắp kết thúc} (Cron thứ Hai 9:00): gửi cho các lịch chuyển tiền sắp hết hạn trong vòng 3 ngày tới.
    \end{enumerate}
    Nhờ kiến trúc này, hệ thống chủ động tương tác với người dùng đúng ngữ cảnh, thay vì chỉ phản ứng khi người dùng mở ứng dụng.
\end{itemize}

\section{Rủi ro đạo đức và giải pháp đảm bảo trách nhiệm AI (Responsible AI)}
Trong hệ thống phần mềm tích hợp AI như Student360, việc xem xét rủi ro đạo đức và thực thi nguyên tắc Responsible AI là bắt buộc. Nhóm phân tích bốn khía cạnh rủi ro cốt lõi:

\begin{enumerate}
    \item \textbf{Thiên lệch trong gợi ý thuật toán (Algorithmic Bias):}
    \begin{itemize}
        \item \textit{Rủi ro:} Thuật toán có xu hướng ưu tiên sinh viên có hồ sơ học thuật xuất sắc, bỏ qua sinh viên hoàn cảnh khó khăn hoặc tích cực hoạt động xã hội.
        \item \textit{Giải pháp:} Hệ thống không cố định công thức chấm điểm. Nhà tài trợ có thể cấu hình trọng số linh hoạt cho từng tiêu chí (nâng trọng số hoàn cảnh khó khăn, hạ trọng số GPA đối với học bổng vượt khó).
    \end{itemize}

    \item \textbf{Khả năng giải thích được (Explainability):}
    \begin{itemize}
        \item \textit{Rủi ro:} Điểm số so khớp (\% tương thích) do AI đưa ra có thể bị coi là ``hộp đen'', khiến sinh viên không hiểu lý do nhận được mức phù hợp thấp.
        \item \textit{Giải pháp:} Hệ thống cung cấp giao diện phân rã điểm số (Match Score Breakdown). Sinh viên xem chi tiết lý do AI đánh giá và AI Agent gợi ý hành động cụ thể để tự cải thiện hồ sơ.
    \end{itemize}

    \item \textbf{Quyền riêng tư dữ liệu và Sự đồng ý (Data Privacy \& User Consent):}
    \begin{itemize}
        \item \textit{Rủi ro:} Dữ liệu học thuật (GPA, bảng điểm) và giao dịch tài chính cá nhân là thông tin cực kỳ nhạy cảm.
        \item \textit{Giải pháp:} Áp dụng nghiêm ngặt nguyên tắc đồng ý của người dùng (User Consent). Dữ liệu cá nhân chỉ được chia sẻ cho nhà tài trợ khi sinh viên chủ động thực hiện hành động nộp hồ sơ (Apply). Sinh viên có quyền kiểm soát, lưu nháp hoặc hủy bỏ hồ sơ bất kỳ lúc nào.
    \end{itemize}

    \item \textbf{Sự phiền nhiễu từ thông báo đẩy (Notification Fatigue):}
    \begin{itemize}
        \item \textit{Rủi ro:} Hệ thống gửi quá nhiều thông báo tự động có thể gây khó chịu và khiến người dùng gỡ ứng dụng.
        \item \textit{Giải pháp:} Triển khai bộ cấu hình tùy chọn thông báo trong Mobile App. Sinh viên có thể tắt/bật hoặc thay đổi tần suất nhận từng loại thông báo cụ thể.
    \end{itemize}
\end{enumerate}

\section{Vấn đề cần giải quyết và định hướng giải pháp}
\subsection{Phân tích các khó khăn vướng mắc}
Trong quá trình xây dựng hệ thống, nhóm phát triển đối mặt với các khó khăn chính sau:
\begin{enumerate}
    \item \textbf{Độ trễ và chi phí tích hợp AI:} Việc gọi trực tiếp các API của mô hình ngôn ngữ lớn (như Gemini) có độ trễ phản hồi tương đối cao và tốn kém tài nguyên.
    \item \textbf{Tính chính xác trong phân loại giao dịch:} Ngôn ngữ tự nhiên mà người dùng nhập vào khi ghi chép chi tiêu thường mang tính viết tắt, không chuẩn hóa ngữ pháp (ví dụ: "an trua 50k", "cf 30k").
    \item \textbf{Đồng bộ trạng thái giao diện thời gian thực:} Giao diện điện thoại di động yêu cầu cập nhật tức thì số dư của các hũ và biểu đồ báo cáo ngay sau khi người dùng thực hiện giao dịch mà không làm chậm trải nghiệm ứng dụng.
    \item \textbf{Bảo mật dữ liệu và phân quyền hệ thống:} Dữ liệu tài chính cá nhân và hồ sơ học thuật/học bổng của sinh viên là các thông tin nhạy cảm. Hệ thống cần đảm bảo an toàn truy cập dữ liệu, chỉ cho phép các đối tượng thích hợp tiếp cận các tập tài nguyên tương ứng (Sinh viên truy cập app cá nhân, Nhà tài trợ quản lý hồ sơ ứng tuyển, Admin quản trị hệ thống).
    \item \textbf{Đảm bảo tính nhất quán và toàn vẹn của dữ liệu giao dịch tài chính:} Khi thực hiện các tác vụ phức tạp như phân bổ thu nhập (income distribution) vào 6 hũ hoặc thực hiện chuyển khoản nội bộ giữa các hũ, việc phát sinh lỗi giữa chừng có thể dẫn đến sai lệch số dư. Ngoài ra, các yêu cầu đầu vào từ người dùng cũng cần được kiểm tra chặt chẽ: ví dụ từ chối giao dịch có thời điểm trong tương lai, hoặc phân bổ thu nhập có tổng các khoản phân bổ lệch quá 0,01 VND so với tổng thu nhập cần chia.
    \item \textbf{Đồng bộ dữ liệu học thuật và bảng điểm cá nhân:} Danh mục trường, khoa/chương trình đào tạo và môn học cần nhất quán giữa Webadmin và Mobile App. Đồng thời, điểm sinh viên nhập phải được kiểm tra theo thang điểm của từng môn, tránh sai lệch GPA và tổng tín chỉ.
\end{enumerate}

\subsection{Giải pháp đề xuất}
Để vượt qua các khó khăn trên, hệ thống đề xuất các giải pháp kỹ thuật ở mức khái niệm sau đây; chi tiết thiết kế và cơ chế hiện thực cụ thể được trình bày đầy đủ tại Chương 3 và Chương 5:
\begin{itemize}
    \item \textbf{Kiến trúc AI Agent kết hợp lưu trữ nhận định tài chính:} FastAPI AI Service đóng vai trò điều phối (orchestration) độc lập, sử dụng các Tool-calling Agent của LangChain để tối ưu hóa việc gọi mô hình, đồng thời lưu bền vững kết quả nhận định tài chính (AI Insights) theo từng chu kỳ báo cáo để tránh phải gọi lại LLM không cần thiết.
    \item \textbf{Phân loại giao dịch bằng LLM kết hợp Rule Engine:} Sử dụng các bộ quy tắc gán cứng (Rule-based) kết hợp kỹ thuật Few-shot Prompting để hướng dẫn mô hình Gemini phân loại các thuật ngữ viết tắt của người dùng một cách chính xác vào 6 lọ tài chính.
    \item \textbf{Đồng bộ giao diện thời gian thực:} Áp dụng cơ chế Invalidation của thư viện quản lý state phía Frontend để tự động làm mới dữ liệu hiển thị ngay khi có thay đổi từ phía máy chủ, không cần tải lại toàn bộ ứng dụng.
    \item \textbf{Xác thực và phân quyền nhiều tầng:} Sử dụng JSON Web Token (JWT) để xác thực người dùng trên các API, kết hợp cơ chế phân quyền theo vai trò (RBAC) để tách bạch phạm vi truy cập dữ liệu giữa sinh viên, nhà tài trợ và quản trị viên.
    \item \textbf{Đảm bảo toàn vẹn dữ liệu giao dịch tài chính:} Áp dụng Database Transaction cho các thao tác chuyển khoản giữa các hũ để đảm bảo tính nguyên tử (Atomicity), đồng thời tính toán lại số dư từ tổng hợp lịch sử giao dịch thực tế thay vì lưu số dư tĩnh, giúp hạn chế sai lệch dữ liệu khi có thao tác điều chỉnh.
    \item \textbf{Thiết kế cơ sở dữ liệu quan hệ chuẩn hóa cao:} Xây dựng cấu trúc dữ liệu với các ràng buộc khóa ngoại (Foreign Keys) chặt chẽ để đồng bộ dữ liệu học thuật phân cấp từ trường học tới sinh viên, đảm bảo tính nhất quán toàn vẹn dữ liệu.
\end{itemize}

\section{Các chức năng chính của hệ thống}
Bốn nhóm chức năng cốt lõi đã nêu tại Mục tiêu và Yêu cầu chức năng ở Chương 1 được hiện thực hóa thành bốn phân hệ tương ứng của hệ thống: Quản lý tài chính cá nhân, Trợ lý tài chính ảo AI, Quản lý học bổng trực tuyến và Quản lý học thuật/Bảng điểm. Danh sách tác nhân tham gia, sơ đồ phân rã chức năng và đặc tả chi tiết từng nghiệp vụ của bốn phân hệ này được trình bày tại Chương 3.
